% !Mode:: "TeX:UTF-8"
%%%论文封面
%%  教育硕士专业学位是按   华南师范大学教师教育学部 2024 年 6 月 6 日颁布的华南师范大学研究生学位论文撰写规范  所制定
%%  学术硕士专业学位是按 华南师范大学数学科学学院 2024 年 5 月 23 日颁布年颁布的 研究生学位论文撰写规范（适用学科：0701数学） 所制定

%%%%%%%%%%%%%%%%%%%%%%%%%%%%%%
 % 专业学位     中文封面内容  
%% “论文分类号”按《中国图书资料分类法》的分类号填写
%%% 具体可查询 http://ztflh.xhma.com
% 分类号 学术型:
% 计算数学 O24   应用数学 O29  
% 概率论与数理统计 O21   运筹学  O22
% 统计学 C8   课程与教学论  G427

\classification{}  % 分类号 
\UDC{} % UDC: 表示国际十进分类法(Universal Decimal Classification),是世界一般报告或者论文的封面在左上角注明分类号，便于信息交换和处理。
%%%  专业学位的 分类号、UDC：不填写
%%%  学术学位的  分类号: 有  计算数学 O24   应用数学 O29   概率论与数理统计 O21    运筹学  O22 统计学 C8   课程与教学论  G427
\confidential{公开}
%% “密级：学位论文密级分为公开、内部、涉密。内部论文和涉密论文严格按照学校审批的保密级别和保密期限填写，如内部*年、秘密*年；其他论文均填公开
\studentID{2020XXXXXX} 
%%   填写学位申请人的学号

 
\type{硕\ 士}      % 可选 硕士 或  博士
%\type{博\ 士}   
 
%%% ==================== 
%%% ====== 中文标题 ======
%%% ==================== 
%%% 论文标题: 应准确概括整篇论文的核心内容，简明扼要，一般不超过25个汉字。必要时可加副标题。
\title{\CASTunderline[\textwidth]{基于 Newman 错误分析的高中生数学解题}
\\ \vspace{0.18cm}
\CASTunderline[\textwidth]{错误研究 ---以三角恒等变换为例}} 
%% 如上是封面中的 论文中文标题, 需要人工分行

\titleab{基于 Newman 错误分析的高中生数学解题错误研究} 
\subtitle{}    % 中文副标题,  排版时自动居右, 这是可选项, 若不需要可注解掉


%%% ==================== 
%%% ====== 英文标题 ======
%%% ==================== 
%%% 英文标题 

\etitle{\CASTunderline[\textwidth]{Title: A Study of Mathematics Problem Solving }
\\ \vspace{0.18cm}
\CASTunderline[\textwidth]{Errors of High School Students Based on Newman}
\\ \vspace{0.18cm}
\CASTunderline[\textwidth]{Error Analysis --- Take the trigonometric constant}
\\ \vspace{0.18cm}
\CASTunderline[\textwidth]{transformations as an example}
} 

 
%%%==================== ==========  
 %%%==========  学位申请人 ==========  
\author{李建国} 
\eauthor{LI Jianguo}
%%% 英（外）文封面作者姓名按姓在前、名在后的格式书写，姓名需写全拼，如LI Jianguo


%%%==================== ==========  
%%%==========  学科专业  ==========  
%%% 学科专业: 申请人所在学科专业，必须严格按参照学科代码/专业学位代码字典中的二级学科名称/专业学位领域名称填写，若所选一级学科/专业学位类别下无二级学科/专业学位领域，则填写一级学科/专业学位类别名称，不可用简称。

%%%    学术学位   
%  \major{计算数学}
% \emajor{Computional Mathematics}
%  \major{应用数学}
% \emajor{Applied Mathematics}
%  \major{基础数学}   
% \emajor{Pure Mathematics}
%  \major{概率论与数理统计}   
% \emajor{Probability Theory and Mathematic Statistics}
%  \major{运筹学与控制论}   
% \emajor{Operations Research and Cybernetics}

%%%       专业学位
\major{教育硕士}   
\emajor{Master of Education}   %Major 仅出现在专业学位中 学科专业 的英文翻译, 时 需要,  学术学位不出现
%  \major{教育博士} % 专业学位名称
%\emajor{Doctor of education of Education}


%%%==================== ==========  
%%%==========  论文研究方向  ==========  
%%% 论文研究方向不能与学科专业名称完全相同，一个研究方向最多12个字，最多两个，以中文分号分隔，总数不超过24个字。
%%%%    学术学位: 研究方向 
%\specialization{微分方程数值解及应用} % 研究方向 
%%%%    专业学位: 
\specialization{学科教学(数学)}  
\especialization{Subject Teaching (Mathematics)} 
 

%%%==================== ==========
 %%%==========   所在学院  ========== 
 %%% 申请人所在二级培养单位的名称，必须填写与学校官网一致的规范名称，不可用简称。
\college{数学科学学院}  
\ecollege{School of Mathematical Sciences}   


%%%==================== ==========
 %%%==========  导师姓名及职称 ========== 
 %%%  填写指导教师的姓名与职称。若有两名指导教师，第一（校内）导师姓名和职称写在前面，第二（或校外）导师姓名和职称写在后面，中间空一格。第一（校内）导师必须与研究生管理系统一致。
\advisor{钟柳强} % 导师姓名
\ProfessionalTitle{教授} % 及职称
\eadvisor{ZHONG Liuqiang}  % 
%英（外）文封面作者姓名按姓在前、名在后的格式书写，姓名需写全拼，如LI Jianguo
\secondadvisor{ } %校外(行业)导师姓名, 仅在专业学位时需要,  学术学位不出现


%%%==================== ==========
 %%%==========  论文完成时间 ==========
 % 填写学位论文最终定稿打印成文的年月，封面的日期用汉字书写，如二〇二四年六月、二〇二四年十二月
\submityear{二〇二五} % 年, 用于中文封面
\submitmonth{五}         % 月, 用于中文封面
\submitengyear{2025}  % 年, 用于英文封面
\submitengmonth{MAY} % 月, 用于英文封面


\maketitle   % 生成中文封面
\makeengtitle   % 生成英文封面